\section{Flat deformation: equations and floor separation}
\label{sec:fl-system-floor}

We first state the asymptotically flat estimate used in the change of
asymptotic model.  The second fundamental form in this estimate need not
have compact support.

\begin{theorem}[Flat deformation estimate]
\label{thm:fl-deformation}
Let \((\Omega_{\mathrm{orig}},q,B)\) be smooth three-dimensional initial
data on a connected orientable manifold with nonempty compact smooth
boundary \(S_0\).  Assume that the metric space including \(S_0\) is
complete and has exactly one asymptotically Euclidean end.  On that end,
in Euclidean coordinates of radius \(r\), suppose that, for some
\(0<\delta<1\),
\begin{equation}
\label{eq:fl-data-decay}
 q-\delta_{\mathrm{Eucl}}=O_\infty(r^{-1}),\qquad
 R_q=O_\infty(r^{-3-\delta}),\qquad
 B=O_\infty(r^{-2}),\qquad \tr_qB=0.
\end{equation}
Here \(O_\infty(r^{-a})\) includes all differentiated symbol estimates:
each Euclidean derivative lowers the indicated order by one.  Extend
\(r\) to a positive smooth function on the whole manifold.  Suppose that
there is \(c_*>0\) such that
\begin{equation}
\label{eq:fl-strict-dec}
 R_q+(\tr_q B)^2-\lvert B\rvert_q^2
 -2\bigl\lvert\operatorname{div}_q
          \bigl(B-(\tr_qB)q\bigr)\bigr\rvert_q
 \ge c_*(1+r)^{-3-\delta},
\end{equation}
and that \(H_q+\tr_{S_0}B<0\) on every component of \(S_0\), with the
normal pointing into the exterior.

For every connected smooth codimension-zero submanifold-with-boundary
\(D\subset\Omega_{\mathrm{orig}}\), closed as a subset and with
\(\operatorname{int}D\subset\operatorname{int}\Omega_{\mathrm{orig}}\),
that contains the entire sufficiently distant end and whose
\emph{full compact intrinsic boundary} is smoothly embedded and
two-sided, use that entire boundary as an enclosing cut.
Allow coincidence with components of \(S_0\), and allow
\(D=\Omega_{\mathrm{orig}}\), whose cut is \(S_0\).  Set
\[
 A_*=\inf_D\Area_q(\partial D).
\]
Then the metric ADM energy satisfies
\[
 E_q\ge \sqrt{\frac{A_*}{16\pi}}.
\]
The boundary, the intermediate cuts, and the compact interior topology
are otherwise unrestricted.
\end{theorem}

The energy here is
\[
 E_q=\frac1{16\pi}\lim_{R\to\infty}
 \int_{\{r=R\}}
   (\partial_jq_{ij}-\partial_iq_{jj})\nu_{\mathrm{Eucl}}^i
   \,\dd A_{\mathrm{Eucl}} .
\]
Its existence follows from the scalar-curvature divergence identity:
the linearized scalar curvature differs from \(R_q\) by an integrable
\(O(r^{-4})\) error under Equation~\eqref{eq:fl-data-decay}.
The proof of Theorem~\ref{thm:fl-deformation} occupies this section and
the subsequent analytic completion.  It is completed in
Theorem~\ref{thm:fc-mass-completion}.

For its proof write \(g=q\) and \(K=B\).  In this flat section the
constraint densities contain no cosmological term:
\begin{equation}
\label{eq:fl-flat-constraints}
 16\pi\mu=R_g+(\tr_gK)^2-\lvert K\rvert_g^2,\qquad
 8\pi J=\operatorname{div}_g\bigl(K-(\tr_gK)g\bigr),\qquad
 M_d=8\pi(\mu-\lvert J\rvert_g).
\end{equation}
In particular \(M_d\ge(c_*/2)(1+r)^{-3-\delta}\) and
\(\tr_gK\) has compact support.  All derivatives and contractions below
refer to \(g\), unless another metric is specified.

\subsection{Full bounding regions and the two thresholds}
\label{subsec:fl-regions}

For a symmetric tensor \(Q\), a two-sided surface \(\Sigma\), and a
specified normal \(n\), put
\[
 \Theta_Q(\Sigma,n)=H_\Sigma(n)+\tr_\Sigma Q.
\]
We will repeatedly use the two identities
\begin{equation}
\label{eq:fl-expansion-signs}
 \Theta_{Q-(c/2)g}=\Theta_Q-c,\qquad
 \Theta_{-Q}(\Sigma,-n)=-\Theta_Q(\Sigma,n).
\end{equation}
The factor \(1/2\) is the reciprocal of the surface dimension.

We use the following barrier and total-region theorem in its
bounding-domain formulation.
We record the boundary conventions to distinguish it from an assertion
about arbitrary, possibly nonbounding, marginal surfaces.

\begin{lemma}[MOTS barriers and the total bounding trapped region]
\label{lem:fl-mots-barriers}
Let \((M,g,Q)\) be a smooth compact three-manifold with boundary
\(\Gamma_-\sqcup\Gamma_+\).  Every connected component of \(M\)
is required to meet \(\Gamma_+\).  The outer part \(\Gamma_+\) has
\(\Theta_Q>0\) with its normal out of \(M\).  The required inner part
\(\Gamma_-\) may be empty; when present it has \(\Theta_Q<0\) with its
normal into \(M\).

Consider all smooth bounding domains that contain a relative collar of
\(\Gamma_-\), avoid \(\Gamma_+\), and whose free boundary has
\(\Theta_Q\le0\) with the normal out of the domain.  If this collection
is nonempty, the closure of its union has smooth compact embedded stable
MOTS free frontier.  It is itself the closure of a bounding domain and
contains every domain in the collection.  If \(\Gamma_-\ne\varnothing\),
the strict inner and outer barriers give a nonempty smooth enclosing
MOTS.  Boundary parts and frontiers may be disconnected.  No constraint
or energy condition on \(Q\) is required.
\end{lemma}

\begin{proof}
This is the total-region theorem of
\cite[Definitions 7.1--7.2 and Theorem 7.3]{anderssonmetzger2009},
together with its strict-barrier existence theorem and stability
conclusion \cite[Theorems 3.1 and 4.1]{anderssonmetzger2009}.
The conventions in Section 2 of that paper allow an empty required
inner boundary for the total-region theorem.  Apply these results on
each connected component of \(M\): its outer part is nonempty, and
strict-barrier existence is used only when its own inner part is
nonempty.  A component with empty inner part contributes a frontier
only if its collection of trapped domains is nonempty.  Take the
finite union of the resulting regions.  Each free frontier is
oriented toward the designated outer part.  Their domains include the
entire required inner boundary, which is not discarded as an empty
relative frontier.
\end{proof}

A component of the complement of a smooth bounding domain that meets
no designated outer face can be filled.  This deletes whole free
boundary components and preserves the expansion bound on all remaining
components.  Thus a full maximal region is filled in this sense.

Choose \(R_0\) so large that every coordinate sphere with \(r\ge R_0\)
has positive expansion for both \(K\) and \(-K\).  Indeed
\[
 H_{\{r=R\}}=\frac2R+O(R^{-2}),\qquad \lvert K\rvert=O(R^{-2}).
\]
No bounding domain with \(\Theta_{\pm K}\le c\le0\) can reach this part
of the end.  At a maximum of its boundary radius, the domain lies inside
the tangent coordinate sphere and its outward normal is radial.
The second-derivative comparison gives mean curvature at least that of
the sphere, contradicting its nonpositive expansion.  All
negative-threshold regions therefore lie in a fixed compact radial
range, independent of threshold and of a sufficiently large truncation.

For
\[
 \max_{S_0}\Theta_K(S_0)<c<0
\]
let \(\mathcal B_c\) be the closed total bounding region for
\(\Theta_K\le c\), requiring all of \(S_0\) on its inner side.
A short relative collar of \(S_0\) is a strict seed.  Apply
Lemma~\ref{lem:fl-mots-barriers} to \(K-(c/2)g\) on a large truncation.
It gives a smooth stable frontier with expansion \(c\).
The regions \(\mathcal B_c\) increase with \(c\).

Fix a threshold \(c_b<0\), to be selected below, and write
\(\mathcal B=\mathcal B_{c_b}\).
In its complement use the entire reversed black frontier and a distant
coordinate sphere as outer barrier faces.  For the tensor \(-K\) the
reversed black expansion is \(-c_b>0\).  There is no required inner face.
The original truncation is connected, and the black region is filled
relative to the distant sphere.  Its remaining complement is therefore
the component containing that sphere; in particular it has a designated
outer face.
Let \(\mathcal W_c\) be the closed total region for
\(\Theta_{-K}\le c<0\) in this complement, and set it equal to the empty
set if its defining collection is empty.  For the shifted tensor
\(-K-(c/2)g\), the reversed black expansion is \(-c_b-c>0\);
hence Lemma~\ref{lem:fl-mots-barriers} applies.  A nonempty
\(\mathcal W_c\) has smooth stable frontier of expansion \(c\), disjoint
from the black faces.  This family is increasing with \(c\), while
\(\mathcal B\) remains fixed.

\begin{lemma}[Right-continuous thresholds]
\label{lem:fl-right-continuity}
Let \(E_c\) be an increasing family of closed subsets of a fixed compact
metric space.  Except at a countable set of parameters,
\[
 E_{c_j}\longrightarrow E_c
 \quad\hbox{in Hausdorff distance whenever }c_j\downarrow c.
\]
At an empty \(E_c\), this means that \(E_{c'}\) is empty for all
sufficiently close \(c'>c\).
\end{lemma}

\begin{proof}
Choose a countable dense set \(x_i\) and \(L\) larger than the diameter
of the compact space.  Set
\[
 d_c(x)=\min\{L,\operatorname{dist}(x,E_c)\},
 \qquad d_c\equiv L\quad\hbox{if }E_c=\varnothing .
\]
For every \(i\), the bounded nonincreasing function \(c\mapsto d_c(x_i)\)
has at most countably many right discontinuities.  Avoid their
countable union.  The functions \(d_c\) are uniformly \(1\)-Lipschitz,
so convergence at the dense set implies uniform convergence by a finite
net argument.  Nesting gives one Hausdorff inclusion; if points in
\(E_{c_j}\) stayed a positive distance from \(E_c\), evaluating the
uniformly convergent distance functions at those points would contradict
\(d_{c_j}=0\) there.  If \(E_c\) is empty, uniform convergence to \(L\)
is incompatible with a zero of any \(d_{c_j}\).
\end{proof}

Choose \(c_b\) at such a continuity point of the black family, then
choose \(c_w<0\) at such a continuity point of the white family for this
fixed black region.  Both choices can be made arbitrarily close to zero.
Write \(\mathcal W=\mathcal W_{c_w}\).  Let \(\Omega\) be the closure
of the component of
\(\Omega_{\mathrm{orig}}\setminus(\mathcal B\cup\mathcal W)\)
that reaches infinity.  Its compact boundary is a finite disjoint union
of whole frontier components:
\begin{equation}
\label{eq:fl-face-data}
 B=\partial\Omega=B_b\sqcup B_w,\qquad
 H+P_B=c_b\ \hbox{on }B_b,\qquad
 H-P_B=c_w\ \hbox{on }B_w,\qquad P_B=\tr_BK .
\end{equation}
Here and below \(\nu\) points into \(\Omega\), and \(H=H_B(\nu)\).
Some black frontier components may be cut off from infinity by the
white region; only the whole components bounding \(\Omega\) are retained.
The white type, or either retained type, may be absent.

Every full enclosing cut of \(\Omega\) is also a full enclosing cut of
\(\Omega_{\mathrm{orig}}\).  More explicitly, its connected closed
exterior submanifold is a submanifold of the original exterior, contains
the distant end, and has the same entire intrinsic boundary.  The
discarded side contains the original \(S_0\) and a collar of it.
Consequently
\begin{equation}
\label{eq:fl-cut-inclusion}
 \Area_g(\Gamma)\ge A_*
 \quad\hbox{for every admissible enclosing cut }\Gamma\hbox{ in }\Omega .
\end{equation}
No minimizing property of the new boundary is asserted or needed.

\begin{lemma}[Outward collars of a maximal frontier]
\label{lem:fl-frontier-collars}
If \(\Sigma\) is a connected component of the smooth frontier of a full
maximal region at threshold \(c\) for \(Q\), it has a smooth outward
foliation \(\Sigma_s\), \(0\le s<s_0\), such that
\[
 \Theta_Q(\Sigma_0)=c,\qquad
 \Theta_Q(\Sigma_s)>c\quad(0<s<s_0).
\]
The leaf parameter is a smooth defining function, and
\(\lvert\dd s\rvert\) is bounded above and away from zero.
\end{lemma}

\begin{proof}
For normal graphs \(v\) over \(\Sigma\), let
\(\Phi(v)=\Theta_Q(\Sigma(v))-c\), pulled back to \(\Sigma\).
Its linearization \(L\) is the scalar MOTS stability operator for
\(Q-(c/2)g\), with principal part \(-\Delta_\Sigma\).
Stability gives a principal eigenvalue \(\lambda\ge0\) and a positive
smooth eigenfunction \(\varphi\).  If \(\lambda>0\), the graphs
\(v=s\varphi\) have
\(\Phi(s\varphi)=s\lambda\varphi+O(s^2)>0\) and positive velocity
for sufficiently small \(s>0\).

If \(\lambda=0\), the kernel of \(L\) is spanned by \(\varphi>0\),
and its adjoint kernel is spanned by \(\varphi^*>0\).
These scalar elliptic principal-eigenfunction properties also follow
by applying the maximum principle after dividing a kernel element by
\(\varphi\), together with Fredholm index zero.  The derivative of the
augmented map
\[
 (v,a)\longmapsto
 \left(\Phi(v)-a,\int_\Sigma v\,\dd A\right)
\]
is
\[
 (v,a)\longmapsto
 \left(Lv-a,\int_\Sigma v\,\dd A\right).
\]
It is an isomorphism.  For a prescribed first component \(f\), pairing
with \(\varphi^*\) uniquely fixes
\(a=-\int\varphi^*f/\int\varphi^*\); adding a unique multiple of
\(\varphi\) then fixes the second component.  The implicit-function
theorem supplies \(v(s),a(s)\) with output \((0,s)\).
Differentiation gives \(a'(0)=0\) and
\(v'(0)=\varphi/\int_\Sigma\varphi>0\).
Thus the graphs are ordered and have constant expansion \(c+a(s)\).
If \(a(s)\le0\) for any sufficiently small \(s>0\), replacing this one
frontier component by that outer graph enlarges the bounding region
while keeping every free component at expansion at most \(c\).
This contradicts maximality.  The claimed bounds on the leaf
parameter follow from positive velocity and compactness.
\end{proof}

Apply this lemma with \(Q=K\) on black faces and \(Q=-K\) on white
faces.  Choose disjoint collars of the finitely many retained components.
Their outward leaf normal is \(\nu_s=\nabla s/\lvert\nabla s\rvert\).

\subsection{Exterior supports and smooth enclosures}
\label{subsec:fl-supports}

\begin{definition}[Exterior-support expansion]
\label{def:fl-exterior-support}
For a closed set \(E\), a smooth exterior support at \(x\in\partial E\)
is a smooth function \(\phi\) near \(x\) satisfying
\(\phi(x)=0\), \(\dd\phi(x)\ne0\), and \(E\subset\{\phi\le0\}\)
locally.  With \(n=\nabla\phi/\lvert\nabla\phi\rvert\), put
\[
 \mathcal E_Q[\phi]
 =\frac{(g^{ij}-n^in^j)\nabla_i\nabla_j\phi}
         {\lvert\nabla\phi\rvert}
   +\tr_gQ-Q(n,n).
\]
The set has exterior-support expansion at most \(c\) if
\(\mathcal E_Q[\phi](x)\le c\) for \emph{every} such support at every
frontier point \(x\).
\end{definition}

For a smooth domain this is precisely an upper bound on its outward
expansion.  A smooth exterior supporting graph has mean curvature no
larger than that of the enclosed graph.  Increasing smooth changes of
defining function preserve \(\mathcal E_Q\), since the extra axial
Hessian has zero tangential trace.

\begin{lemma}[Distance neighborhoods with an additive support error]
\label{lem:fl-support-parallel}
Suppose that the frontier of a closed set \(E\) lies in a fixed compact
interior region of smooth data \((M,g,Q)\), and has exterior-support
expansion at most \(c\).  For sufficiently small \(z>0\), the distance
neighborhood \(E_z=\{\operatorname{dist}(\,\cdot\,,E)\le z\}\) has
exterior-support expansion at most \(c+Cz\).
The constant and the allowable radius depend only on the compact
ambient geometry and \(Q\), independently of the support's curvature
and the regularity of \(E\).
\end{lemma}

\begin{proof}
Let \(\phi\) support \(E_z\) at \(x'\), and choose a nearest point
\(x\in E\), joined to \(x'\) by a length-\(z\) minimizing unit geodesic
\(\gamma\).  Choose \(z\) below the injectivity radius of the fixed
compact neighborhood.  Since the ball of radius \(z\) centered at
\(x\) is inside \(E_z\), the support normal at \(x'\) is
\(\dot\gamma(z)\).  Let \(P_t\) denote parallel transport along
\(\gamma\).

For \(X\in T_xM\), solve the Jacobi boundary-value problem
\[
 \nabla_{\dot\gamma}^2J_X+
       \operatorname{Rm}(J_X,\dot\gamma)\dot\gamma=0,\qquad
 J_X(0)=X,\qquad J_X(z)=P_zX.
\]
In a parallel frame its integral equation is
\[
 j_X(t)=X+tA_X-\int_0^t(t-s)\mathcal R(s)j_X(s)\,\dd s,\qquad
 A_X=\frac1z\int_0^z(z-s)\mathcal R(s)j_X(s)\,\dd s.
\]
For small \(z\), absorption gives
\(\sup_t\lvert j_X(t)\rvert\le2\lvert X\rvert\) and
\(\lvert A_X\rvert\le Cz\lvert X\rvert\).
Moreover \(\langle J_X,\dot\gamma\rangle\) is affine with equal endpoint
values.  Hence \(A_X\) is perpendicular to \(\dot\gamma(0)\).
We can prescribe the first jet of a local unit vector field \(V\) by
\[
 V(x)=\dot\gamma(0),\qquad \nabla_XV(x)=A_X .
\]
To realize this jet, use a frame obtained by radial parallel transport
from \(x\), choose affine coefficients with the prescribed first jet,
and normalize their length.  The perpendicularity just proved means
that normalization preserves the first derivative.  The second
derivatives of this unit field have a uniform bound on a small ball.

For \(F_z(y)=\exp_y(zV(y))\), the Jacobi construction gives
\begin{equation}
\label{eq:fl-support-map}
 F_z(x)=x',\qquad (\dd F_z)_x=P_z,\qquad
 \lvert(\nabla\dd F_z)_x\rvert\le Cz .
\end{equation}
For the last estimate, differentiate the exponential map twice,
including the bounded jets of \(V\).  At \(z=0\) the map is the identity
and its covariant second derivative is zero; the first derivative in
\(z\) is uniformly bounded on these jets.  This construction never
uses derivatives of the supporting function.

Every \(y\in E\) near \(x\) is moved a distance \(z\), so
\(F_z(y)\in E_z\).  Thus \(\psi=\phi\circ F_z\) supports \(E\) at \(x\).
Equation~\eqref{eq:fl-support-map} makes its gradient length exactly
that of \(\phi\) at \(x'\), and its normal is \(\dot\gamma(0)\).
The covariant chain rule reads
\[
 \Hess\psi(X,Y)=\Hess\phi(P_zX,P_zY)
       +\dd\phi\bigl((\nabla\dd F_z)(X,Y)\bigr).
\]
Trace over \(\dot\gamma(0)^\perp\) and divide by the common gradient
length.  The first term is exactly the normalized tangential Hessian
of \(\phi\); the second costs \(Cz\).  Parallel transport changes the
tangential trace of \(Q\) by at most \(Cz\).  Therefore
\(\lvert\mathcal E_Q[\psi](x)-\mathcal E_Q[\phi](x')\rvert\le Cz\).
The assumed bound for \(\psi\) proves the result.  In particular there
is no error proportional to \(z\lvert\Hess\phi\rvert\).
\end{proof}

\begin{lemma}[Smooth enclosure from all exterior supports]
\label{lem:fl-support-enclosure}
Let \(M\) be a smooth compact connected three-manifold with boundary
\(\Gamma_-\sqcup\Gamma_+\), where \(\Gamma_+\ne\varnothing\).
Let \(E\ne\varnothing\) be closed, contain a relative collar of the
required \(\Gamma_-\) if it is present, and have positive distance from
\(\Gamma_+\).  Suppose every exterior support of its interior frontier
has \(Q\)-expansion at most \(c\).  If
\[
 c'>c,\qquad
 \Theta_Q(\Gamma_+)>c'
\]
with the outward normal of \(M\), then \(E\) is strictly enclosed by a
smooth bounding domain whose free frontier has expansion \(c'\).
In particular it lies in the full closed trapped region at threshold
\(c'\).  The enclosing domain and frontier may be disconnected.
\end{lemma}

\begin{proof}
Take \(u>0\) sufficiently small that all new distance frontiers are in
a fixed compact interior neighborhood, avoid \(\Gamma_+\), and satisfy
\(c+Cu<c'\) in Lemma~\ref{lem:fl-support-parallel}.
The collar of \(\Gamma_-\) ensures that a nearest segment ending on a
new distance frontier starts on the interior frontier of \(E\);
no extension across the required boundary is needed.

Uniformly approximate \(\operatorname{dist}(\,\cdot\,,E)\) by a smooth
function with error below \(u/32\), using a finite coordinate cover,
mollification, and a partition of unity.  A regular value between
\(u/3\) and \(u/2\) gives a smooth relative neighborhood \(U\) with
\[
 E_{u/4}\subset U\subset E_{3u/4}.
\]
Fill any component of its complement that meets no designated outer
face.  This may enlarge its interior beyond \(E_{3u/4}\), but only
deletes free frontier components; all remaining frontier points still
have distance between \(u/4\) and \(3u/4\) from \(E\).

Choose a smooth \(\eta\), equal to one on that compact free frontier,
with \(0\le\eta\le1\) and support in
\(\{\operatorname{dist}(\,\cdot\,,E)<u\}\).
For a sufficiently large finite \(A>0\), set
\[
 Q_A=Q-\tfrac12(c'+A\eta)g .
 \]
The free frontier of \(U\) is strictly negative for \(Q_A\), while
the outer faces remain strictly positive.  Apply the strict-barrier
part of Lemma~\ref{lem:fl-mots-barriers} to each component of the
complement of \(U\).  Each meets an outer face by filling and has
nonempty inner boundary by connectedness of \(M\) and \(U\ne\varnothing\).
The resulting smooth frontier \(\Sigma\) encloses \(U\) and satisfies
\begin{equation}
\label{eq:fl-modified-mots}
 \Theta_Q(\Sigma)=c'+A\eta\ge c'.
\end{equation}

If \(\Sigma\) enters the support of the modification, let
\(z=\min_\Sigma\operatorname{dist}(\,\cdot\,,E)<u\).
We have \(z\ge u/4>0\), and \(E_z\) lies on the enclosed side of
\(\Sigma\).  Indeed, a point at distance less than \(z\) from \(E\)
can be joined to \(E\) by a path of length less than \(z\); crossing
\(\Sigma\) would give a point of \(\Sigma\) at smaller distance.
At a minimizing point, a defining function of \(\Sigma\) is therefore
an exterior support of \(E_z\).  It has expansion at most
\(c+Cz<c'\), contradicting Equation~\eqref{eq:fl-modified-mots}.
Thus \(\eta=0\) on all of \(\Sigma\), giving the desired expansion.
The inclusion \(E_{u/4}\subset U\) makes the enclosure strict.
Filling remaining components without an outer face again only deletes
whole smooth boundary components.
\end{proof}

This lemma applies after truncating an exterior by a strict distant
sphere.  For the white construction the reversed black faces belong to
\(\Gamma_+\), and there is no required \(\Gamma_-\).
An empty weak set requires no enclosure.
Also, filling a closed set preserves any bound on all exterior
supports: its new frontier is contained in the old closed set, and a
support of the larger filled set also supports the old one at every
remaining frontier point.

\subsection{Small offsets and the prepared data}
\label{subsec:fl-offsets}

\begin{lemma}[Threshold preparation and offsets]
\label{lem:fl-threshold-preparation}
The thresholds and the prepared exterior above can be chosen so that
both full threshold families are right-continuous at their chosen
thresholds, the retained faces have disjoint collars as in
Lemma~\ref{lem:fl-frontier-collars}, and Equation~\eqref{eq:fl-cut-inclusion}
holds.  There are a smooth compactly supported \(C\), constants
\(C_b>0>C_w\), positive collar constants \(k_B\), and a smooth
\(\rho>0\), equal far out to a positive constant times \(r^{-3-\delta}\),
such that
\begin{align}
\label{eq:fl-offset-shape}
 &C_w\le C\le C_b,\qquad
 C=C_b-k_Bs\ \hbox{near }B_b,\qquad
 C=C_w+k_Bs\ \hbox{near }B_w,\\
\label{eq:fl-face-heights}
 &b_-=c_b-C_b<0,\qquad b_+=-c_w-C_w>0,\\
\label{eq:fl-margin}
 &\lvert(h+C)\tr K\rvert+\lvert\dd C\rvert+\rho
       \le M_d\qquad\hbox{for all }b_-\le h\le b_+ .
\end{align}
All these data are fixed before the deformation parameters.
\end{lemma}

\begin{proof}
Fix a compact set containing the common radial range of all the
negative-threshold regions and \(\supp(\tr K)\).
On this compact set \(M_d\) has a positive minimum.  Restrict
\(c_b,c_w\) sufficiently close to zero so that
\(\max\{\lvert c_b\rvert,\lvert c_w\rvert\}\lvert\tr K\rvert\)
uses less than one quarter of this minimum.  The two successive
right-continuity choices remain possible by
Lemma~\ref{lem:fl-right-continuity}.
Fix the regions and their collars.

On each collar take a smooth cutoff \(\kappa(s)\), equal to one near
the face and zero near its far end, and choose \(k>0\) such that
\(1-ks>0\) throughout that collar.  A black contribution
\(a_0\kappa(s)(1-ks)\) and a white contribution
\(-a_0\kappa(s)(1-ks)\), extended by zero and summed over disjoint
collars, give the required \(C\) with \(C_b=a_0\), \(C_w=-a_0\),
and \(k_B=a_0k\).  The \(C^1\) norm tends to zero with \(a_0\), after
the collars have been fixed.  Since
\[
 \lvert h+C\rvert
 \le\max\{\lvert c_b\rvert,\lvert c_w\rvert\}+C_b-C_w ,
\]
choose \(a_0\) small enough for the first two terms of
Equation~\eqref{eq:fl-margin} to be at most \(M_d/2\).
They vanish outside a compact set.  Finally multiply a fixed smooth
positive symbol weight, equal to \(r^{-3-\delta}\) on the end, by a
small constant so that \(\rho\le M_d/2\) everywhere.
\end{proof}

When a retained type is empty its constants still serve as bounds and
its collar terms are omitted.  On the faces, the exact identities are
\begin{equation}
\label{eq:fl-compatible-values}
 b_-+C=P_B+H\quad\hbox{on }B_b,\qquad
 b_++C=P_B-H\quad\hbox{on }B_w .
\end{equation}

\subsection{Variables and the physical system}
\label{subsec:fl-equations}

Extend the end radius smoothly with \(r\ge1\), and write \(\Omega_R\)
for the truncation with outer boundary \(S_R\).
First fix \(0<\epsilon<1\), then choose an integer
\(N\ge4\), \(N>2/\epsilon\), and finally take
\(R\ge R_{\min}(N,\epsilon)\) as required by the estimates.
A \emph{polynomial constant} means a quantity bounded by
\(C(1+N)^m\), with \(C,m\) depending on the fixed prepared geometry,
profiles and \(\epsilon\), but not on \(R\), the solution, or the
homotopy parameter.  Later estimates allowed arbitrary dependence on
fixed \(N\) will be identified separately.

Take a smooth nonincreasing \(\vartheta:\mathbb R\to[0,1]\), equal
to one on \((-\infty,0]\) and zero on \([1,\infty)\), and put
\begin{equation}
\label{eq:fl-profile}
 p(t)=N\vartheta(Nt),\qquad
 l(t)=\exp\!\left(\int_0^t p(z)\,\dd z\right),\qquad
 L(t)=e^{2t}l(t),\qquad
 \ell=e^{-\epsilon N},\qquad \tau=\ell^{3/2}.
\end{equation}
Thus \(l=e^{Nt}\) for \(t\le0\), \(l\) is constant for \(t\ge1/N\),
and every fixed-order derivative of \(p\) is polynomially bounded.
On \(t\ge-\epsilon\), one has \(\ell\le l\le e\).

We solve for \((f,Z)\); the formulas below define \(t\) implicitly and
then define the graph metric, flux, and source.
For functions \(f,t\) define
\begin{align}
\label{eq:fl-variables}
 &\sigma=\lvert\dd f\rvert,\qquad D=1+l^2\sigma^2,\qquad
 d=D^{-1},\notag\\
 &
 a=\frac{l\sigma}{\sqrt D},\quad v=a^2=1-d,\quad
 w=\frac{l\nabla f}{\sqrt D},\qquad \chi=\frac{4d}{4+pv},\\
\label{eq:fl-Z-change}
 &Z=t+\tfrac18\log D,\qquad
 h=\tau f,\qquad u=L\sqrt D,\qquad U=l\sqrt D,\\
\label{eq:fl-graph-data}
 &\bar g=g+l^2\dd f^2,\qquad
 \hat g=e^{4t}\bar g,\qquad
 H^f=\frac{l}{\sqrt D}\Hess f,\qquad S=K+H^f .
\end{align}
For \(\sigma>0\), let \(e=\nabla f/\sigma\) and
\(A_j=I+(j-1)e\otimes e\).  The matrices actually used have
direction-free expressions
\begin{equation}
\label{eq:fl-matrices}
 A_d=I-w\otimes w,\qquad
 A_\chi=I-\frac{p+4}{4+p\lvert w\rvert^2}\,w\otimes w .
\end{equation}
Vectors and covectors are identified using \(g\).
For fixed \(\dd f\),
\[
 \frac{\partial Z}{\partial t}=1+\frac{pv}{4}>0.
\]
The defining function of \(Z\) tends to the corresponding infinite
endpoint as \(t\to\pm\infty\).  Thus Equation~\eqref{eq:fl-Z-change}
defines a unique smooth \(t=t(Z,\dd f)\) for every finite input,
before imposing a floor.  All the quantities in the actual equations
are smooth at \(\dd f=0\).

The physical equations, in the unknowns \((f,Z)\), are
\begin{align}
\label{eq:fl-trace}
 &\tr_{A_\chi}S=F,\qquad F=h+C,\\
\label{eq:fl-scalar}
 &\operatorname{div}V=u\Xi,\qquad
 V=uA_\chi\bigl(4\nabla Z+K(w,\cdot)\bigr).
\end{align}
The notation \(\tr_A\) means contraction with the contravariant
matrix \(A\).  The trace Hessian coefficient at fixed \(Z\) is
\(l\sqrt d\,A_\chi\).  The scalar principal coefficient at fixed
trace field is \(4uA_\chi\).  Their eigenvalues are positive at all
finite inputs; they are uniformly elliptic on bounded input ranges.
The source below contains quadratic second derivatives of \(f\).
Separate positivity of these scalar principal matrices is therefore
not being used as a coupled-system regularity theorem.

At \(\sigma>0\), decompose \(S\) into
\[
 P=S|_{e^\perp\times e^\perp},\qquad
 M=S(e,\cdot)|_{e^\perp},\qquad q_1=S(e,e),\qquad
 \tr P=F-\chi q_1 .
\]
Write \(x=\partial_e t\), \(y=\nabla_\perp t\), and
\(P^{\mathrm{tf}}=P-\frac12(\tr P)g|_{e^\perp}\).
Use the quadratic form
\begin{align}
\label{eq:fl-quadratic-form}
 \mathcal T={}&
 \tfrac12\lvert P^{\mathrm{tf}}\rvert^2
 +\lvert M+(p+1)ay\rvert^2+[4(p+1)-v]\lvert y\rvert^2 \notag\\
 &+4(p+1)d\left(x+\frac{\chi a q_1}{4d}\right)^2
 +\tfrac34\left(F-\frac{\chi q_1}{3}\right)^2
 +\frac{4d(9-d)}{3(4+pv)^2}q_1^2 .
\end{align}
Here the actual trace \(F=\tr_{A_\chi}S\) is substituted.  Its smooth
extension and polynomial coercivity are proved below.

Fix \(3/4<\gamma<1\) and smooth positive symbol weights
\begin{equation}
\label{eq:fl-weights}
 \rho_0\asymp r^{-3-\delta},\qquad \rho_1\asymp r^{-2\gamma}
\end{equation}
on the end.  Define
\begin{equation}
\label{eq:fl-source}
\begin{aligned}
 &m_0(t)=\vartheta\bigl(N(t+\epsilon)\bigr),\qquad
 W_N=\rho_0+v\rho_1,\\
 &\Xi=\delta_0\mathcal T+\rho+\delta_0\tau a\sigma
                  -m_0(t)\Pi_NW_N .
\end{aligned}
\end{equation}
The constant \(\delta_0>0\) is chosen small independently of \(N,R\);
\(\Pi_N=C_\Pi(1+N)^m\) will then be chosen sufficiently large.
On the allowed penalty band
\(-\epsilon\le t\le-\epsilon+1/N\),
\[
 \ell\le l\le e\ell,\qquad L\asymp_\epsilon\ell .
\]
At \(S_R\) prescribe \(f=Z=0\).
On the inner boundary prescribe \(h=b_-\) on \(B_b\) and \(h=b_+\)
on \(B_w\), and
\begin{equation}
\label{eq:fl-boundary}
 \frac{V_\nu}{u}=\mathcal B_F,\qquad
 \mathcal B_F=-H+w_\nu(F-P_B)-N_Bd,\qquad
 N_B>1+\sup_B\lvert H\rvert .
\end{equation}

\subsection{The graph identity and its coercive quadratic form}
\label{sec:fl-alg}

We supply the pointwise calculations used in the flat deformation.
All derivatives, norms, traces, contractions, and divergences in this
subsection refer to \(g\), unless a different metric is displayed.
A one-form in a divergence is identified with its \(g\)-dual vector.
For symmetric covariant tensors define
\[
 [A,B]=\langle A,B\rangle-(\tr A)(\tr B),
 \qquad p_A=A-(\tr A)g.
\]
In particular, the flat constraints have the normalization
\begin{equation}
 R_g=16\pi\mu+[K,K],
 \qquad \operatorname{div}p_K=8\pi J.
 \label{eq:fl-alg-constraints}
\end{equation}
There is no cosmological term in these identities.

Write \(p=(\log l)'(t)\), so that \(\dd l=pl\,\dd t\), and retain the
graph and trace quantities defined above.
The identities in this subsection require only \(l>0\), \(0\le p\le N\),
and the trace equation \(\tr_{A_\chi}S=F\).
At a point where \(\dd f\ne0\), use an orthonormal frame with first
vector \(e=\nabla f/|\dd f|\), and decompose
\[
 S=\begin{pmatrix}q_1&M^{\mathsf T}\\M&P\end{pmatrix},
 \qquad \dd t=x e^\flat+y,\qquad y(e)=0.
\]
Here \(P\) is a symmetric tensor on the two-dimensional plane
\(e^\perp\), and
\begin{equation}
 v=a^2=1-d,\qquad
 \chi=\frac{4d}{4+pv},\qquad
 P_0:=\tr P=F-\chi q_1.
 \label{eq:fl-alg-blocks}
\end{equation}
We use \(P^{\mathrm{tf}}=P-\frac12(\tr P)g|_{e^\perp}\).
The symbol \(|\dd t|_{A_d}^2\) means
\(\langle A_d\nabla t,\nabla t\rangle=dx^2+|y|^2\).

\begin{lemma}[Differentiation and flux identities]
\label{lem:fl-alg-flux}
For \(U=l\sqrt D\), one has
\begin{align}
 \tfrac12\nabla\log D
   &=H^f(w,\cdot)+pv\,\dd t,
   \label{eq:fl-alg-D-derivative}\\
 \nabla_iw_j
   &=(H^f A_d)_{ij}+pd\,(\dd t)_i w_j,
   \label{eq:fl-alg-w-derivative}\\
 \nabla\log u
   &=H^f(w,\cdot)+(2+p+pv)\dd t.
   \label{eq:fl-alg-u-derivative}
\end{align}
Set
\[
 k=-H^f-p(\dd t\otimes w^\flat+w^\flat\otimes\dd t),
 \qquad Q=K-k.
\]
Then the flux \(V=uA_\chi(4\nabla Z+K(w,\cdot))\) satisfies
\begin{align}
 \frac{V}{u}
    &=p_Q(w,\cdot)+4A_d\nabla t+wF,
    \label{eq:fl-alg-flux-identity}\\
 \frac{\operatorname{div}(uw)}u
    &=F+(1-\chi)q_1-\tr K+2(p+1)ax.
    \label{eq:fl-alg-uw-divergence}
\end{align}
Every expression asserted to be an identity extends smoothly across
\(\dd f=0\).
\end{lemma}

\begin{proof}
Differentiate \(D=1+l^2|\dd f|^2\), use \(\dd l=pl\,\dd t\), and
substitute \(H^f=l\Hess f/\sqrt D\) and
\(w=l\nabla f/\sqrt D\). This gives
\eqref{eq:fl-alg-D-derivative}. Differentiating the last formula for
\(w\), and using \(1-v=d\), gives
\[
 \nabla_iw_j
 =H^f_{ij}-H^f(w,\cdot)_i w_j+pd\,(\dd t)_i w_j.
\]
Since \(A_d=I-w\otimes w^\flat\), this is
\eqref{eq:fl-alg-w-derivative}.
Equation \eqref{eq:fl-alg-u-derivative} follows from
\(u=e^{2t}l\sqrt D\).
Taking the trace of \eqref{eq:fl-alg-w-derivative} and adding
\(\langle\nabla\log u,w\rangle\) gives
\[
 \frac{\operatorname{div}(uw)}u
 =\tr H^f+2(p+1)ax.
\]
Now \(\tr H^f=F+(1-\chi)q_1-\tr K\), proving
\eqref{eq:fl-alg-uw-divergence}.

From \(Z=t+\frac18\log D\), Equation
\eqref{eq:fl-alg-D-derivative} yields
\[
 4\nabla Z+K(w,\cdot)=(4+pv)\nabla t+S(w,\cdot).
\]
Its image under \(A_\chi\) has axial and transverse components
\[
 4dx+\chi aq_1,\qquad (4+pv)y+aM,
\]
because \(\chi(4+pv)=4d\).
On the other hand,
\[
 Q_{ee}=q_1+2pa x,\qquad Q_{e\perp}=M+pa y,
 \qquad Q_\perp=P.
\]
The same two components of
\(p_Q(w,\cdot)+4A_d\nabla t+wF\) are
\(-aP_0+4dx+aF\) and \(a(M+pa y)+4y\).
Equation \eqref{eq:fl-alg-blocks} identifies these with the preceding
display, proving \eqref{eq:fl-alg-flux-identity}.
Smoothness follows either from the original definitions or from
\begin{equation}
 A_\chi
 =I-\frac{p+4}{4+p|w|^2}\,w\otimes w^\flat.
 \label{eq:fl-alg-smooth-matrix}
\end{equation}
In particular no choice of \(e\) is needed at \(w=0\).
\end{proof}

\begin{proposition}[Scalar-curvature identity]
\label{prop:fl-alg-scalar}
Let
\[
 \bar g=g+l^2\dd f^2,\qquad \hat g=e^{4t}\bar g.
\]
For \(\mathcal T\) in \eqref{eq:fl-quadratic-form}, the exact identity
is
\begin{equation}
 \frac12 e^{4t}R_{\hat g}
 =8\pi\bigl(\mu+J(w)\bigr)+\mathcal T+w(F)
       -F\tr K-\frac{\operatorname{div}V}{u}.
 \label{eq:fl-alg-identity}
\end{equation}
\end{proposition}

\begin{proof}
\emph{The stationary calculation.}
On a product with coordinate \(s_0\), introduce the stationary Lorentz
metric
\[
 \mathbf g=-l^2(\dd s_0-\dd f)^2+\bar g
           =-l^2\dd s_0^2+2l^2\dd s_0\dd f+g.
\]
The \(s_0\)-slices have induced metric \(g\), shift
\(\beta=l^2\nabla f=Uw\), lapse \(U=l\sqrt D\), and unit normal
\(n=U^{-1}(\partial_{s_0}-\beta)\). We choose the sign
\(k(X,Y)=\mathbf g(\mathbf\nabla_Xn,Y)\).
Since the slice metric is independent of \(s_0\),
\[
 k=-\frac{\operatorname{sym}\nabla\beta}{U}
   =-H^f-p(\dd t\otimes w^\flat+w^\flat\otimes\dd t).
\]
Here \(\operatorname{sym}T=(T+T^{\mathsf T})/2\).

For completeness, put \(\vartheta=\tr k\). The contracted Gauss
equation and the normal expansion equation, with this sign convention,
read
\[
 R_{\mathbf g}=R_g-2\Ric_{\mathbf g}(n,n)-|k|^2+\vartheta^2,
 \qquad
 \Ric_{\mathbf g}(n,n)=-n(\vartheta)-|k|^2+U^{-1}\Delta_gU.
\]
The lapse term in the second formula follows from the acceleration
\(\mathbf\nabla_n n=\nabla\log U\); its divergence contribution is
\(\Delta_g\log U+|\nabla\log U|^2=U^{-1}\Delta_gU\).
Stationarity gives \(Un(\vartheta)=-\beta(\vartheta)\), while
\(\operatorname{div}\beta=-U\vartheta\).
Combining these equalities gives
\begin{equation}
 UR_{\mathbf g}
 =U(R_g+[k,k])
   -2\operatorname{div}\bigl(\nabla U+\vartheta\beta\bigr).
 \label{eq:fl-alg-stationary-scalar}
\end{equation}

In the coordinate \(s=s_0-f\), the same metric is
\(-l^2\dd s^2+\bar g\). Direct contraction of its warped-product
curvature gives
\(R_{\mathbf g}=R_{\bar g}-2l^{-1}\bar\Delta l\).
The determinant and inverse metric of \(\bar g\) give, for a scalar
\(\varphi\),
\[
 \bar\Delta\varphi
 =D^{-1/2}\operatorname{div}(\sqrt D\,A_d\nabla\varphi).
\]
Moreover, Equations \eqref{eq:fl-alg-D-derivative} and
\eqref{eq:fl-alg-w-derivative} imply
\[
 \nabla U-\frac Ul A_d\nabla l=-k(\beta,\cdot).
\]
For example, after division by \(U\) its left side is
\[
 H^f(w,\cdot)+p\{v\,\dd t+w^\flat\langle w,\nabla t\rangle\}
 =-k(w,\cdot).
\]
Substituting into \eqref{eq:fl-alg-stationary-scalar} therefore yields
\begin{equation}
 R_{\bar g}=R_g+[k,k]+\frac2U\operatorname{div}(p_k\beta).
 \label{eq:fl-alg-graph-scalar}
\end{equation}

\emph{The constraint substitution.}
The symmetric part of \(\nabla\beta\) is \(-Uk\), so
\eqref{eq:fl-alg-constraints} gives
\[
 \operatorname{div}(p_K\beta)
 =8\pi UJ(w)-U[K,k].
\]
Use \(Q=K-k\), \(p_Q=p_K-p_k\), and
\([Q,Q]=[K,K]+[k,k]-2[K,k]\) in
\eqref{eq:fl-alg-graph-scalar}. The result is
\begin{equation}
 R_{\bar g}
 =16\pi(\mu+J(w))+[Q,Q]
      -\frac2U\operatorname{div}(Up_Q(w,\cdot)).
 \label{eq:fl-alg-constraint-graph}
\end{equation}

\emph{The conformal calculation.}
In dimension three,
\[
 \frac12 e^{4t}R_{\hat g}
 =\frac12R_{\bar g}-4\bar\Delta t-4|\dd t|_{A_d}^2.
\]
Replacing the weight \(U\) in the last flux of
\eqref{eq:fl-alg-constraint-graph} by \(u=e^{2t}U\) adds
\(2p_Q(w,\nabla t)\).
Writing \(\bar\Delta t\) in the preceding divergence form, and using
\(\log(u/\sqrt D)=2t+\log l\), gives
\[
 -4\bar\Delta t-4|\dd t|_{A_d}^2
 =-\frac1u\operatorname{div}(4uA_d\nabla t)
        +4(p+1)|\dd t|_{A_d}^2.
\]
By \eqref{eq:fl-alg-flux-identity}, the two fluxes obtained so far
sum to \(V-uwF\). Its divergence contributes
\[
 \frac1u\operatorname{div}(uwF)
 =w(F)+F\frac{\operatorname{div}(uw)}u.
\]
It follows that the quadratic term remaining after the displayed
\(-F\tr K\) in \eqref{eq:fl-alg-identity} is
\begin{equation}
 \frac12[Q,Q]+2p_Q(w,\nabla t)+4(p+1)|\dd t|_{A_d}^2
       +F\{F+(1-\chi)q_1+2(p+1)ax\}.
 \label{eq:fl-alg-invariant-quadratic}
\end{equation}
The next calculation identifies it with \(\mathcal T\), completing
the proof.
\end{proof}

\begin{lemma}[Square completion and uniformity]
\label{lem:fl-alg-coercivity}
The expression \eqref{eq:fl-alg-invariant-quadratic} equals
\(\mathcal T\) in \eqref{eq:fl-quadratic-form}. It is smooth in the
actual variables at \(\dd f=0\). For an absolute constant \(C\), if
\(0\le p\le N\), then
\begin{equation}
 F^2+|P|^2+|M|^2+dq_1^2+|\dd t|_{A_d}^2
       \le C(N+4)^2\mathcal T.
 \label{eq:fl-alg-general-coercivity}
\end{equation}
At a point with \(\dd t=0\), put
\(\mathcal E=|P^{\mathrm{tf}}|^2+|M|^2+F^2+\chi q_1^2\).
There is the stronger bound
\begin{equation}
 \frac12\mathcal E\le\mathcal T\le\mathcal E.
 \label{eq:fl-alg-floor-comparison}
\end{equation}
These constants are independent of \(d\), \(p\), \(N\), \(l\), and
the size of \(\dd f\).
\end{lemma}

\begin{proof}
Using \(P_0=F-\chi q_1\), direct contraction gives
\begin{align*}
 \frac12[Q,Q]
 &=\frac12|P^{\mathrm{tf}}|^2-\frac14P_0^2
   +|M+pa y|^2-P_0(q_1+2pa x),\\
 2p_Q(w,\nabla t)&=-2aP_0x+2aM\cdot y+2pv|y|^2.
\end{align*}
Consequently \eqref{eq:fl-alg-invariant-quadratic} is
\begin{align}
 \mathcal T={}&\frac12|P^{\mathrm{tf}}|^2
     +|M+(p+1)a y|^2+[4(p+1)-v]|y|^2
       +4(p+1)dx^2+2(p+1)a\chi q_1x \notag\\
 &\quad+\frac34F^2-\frac12\chi Fq_1
                +(\chi-\chi^2/4)q_1^2 .
 \label{eq:fl-alg-expanded-quadratic}
\end{align}
Completion of its axial and trace squares is exact because
\[
 \chi-\frac{\chi^2}{4}
 =\frac{(p+1)\chi^2v}{4d}+\frac{\chi^2}{12}
                  +\frac{4d(9-d)}{3(4+pv)^2}.
\]
This proves the asserted formula for \(\mathcal T\).

For smoothness, rewrite
\((1-\chi)q_1=\tr S-F\) in
\eqref{eq:fl-alg-invariant-quadratic}, and
\(ax=\langle w,\nabla t\rangle\).
Every term is then an invariant smooth expression in \(S,w,\dd t,p\),
with \(F=\tr_{A_\chi}S\) and \(A_\chi\) given by
\eqref{eq:fl-alg-smooth-matrix}.
At \(w=0\) its value is
\[
 \mathcal T=\frac12|S|^2+\frac12(\tr S)^2
                               +4(p+1)|\dd t|^2.
\]
This also defines the direction-independent value at a critical point
of \(f\).

To prove \eqref{eq:fl-alg-general-coercivity}, denote the translated
variables in the axial, trace, and mixed squares by
\[
 z=x+\frac{\chi a q_1}{4d},\qquad
 r_F=F-\frac{\chi q_1}{3},\qquad
 m_F=M+(p+1)a y.
\]
The square formula bounds \(|P^{\mathrm{tf}}|^2,|m_F|^2\),
\((p+1)|y|^2\), \((p+1)dz^2\), and \(r_F^2\) by absolute multiples
of \(\mathcal T\). Its last coefficient satisfies
\[
 \frac{4d(9-d)}{3(4+pv)^2}
                    \ge \frac{32d}{3(N+4)^2},
\]
so \(dq_1^2\le C(N+4)^2\mathcal T\).
Since \(\chi\le d\le1\), one has
\[
 dx^2\le 2dz^2+\frac{\chi^2v}{8d}q_1^2
                    \le2dz^2+\frac18dq_1^2,
 \qquad
 F^2\le2r_F^2+\frac29\chi^2q_1^2.
\]
These estimates control \(x,F\), and
\(|P|^2=|P^{\mathrm{tf}}|^2+\frac12(F-\chi q_1)^2\).
Finally,
\[
 |M|^2\le2|m_F|^2+2(p+1)^2v|y|^2
                               \le C(N+1)\mathcal T.
\]
They prove \eqref{eq:fl-alg-general-coercivity}.

If \(\dd t=0\), Equation \eqref{eq:fl-alg-expanded-quadratic} becomes
\[
 \mathcal T=\frac12|P^{\mathrm{tf}}|^2+|M|^2
       +\frac34F^2-\frac12\chi Fq_1
                      +\chi(1-\chi/4)q_1^2.
\]
Using
\(\frac12|\chi Fq_1|\le\frac14F^2+\frac14\chi^2q_1^2\)
and \(0<\chi\le1\) proves both bounds in
\eqref{eq:fl-alg-floor-comparison}.
\end{proof}

\subsection{Pointwise estimates at an interior floor contact}
\label{sec:fl-alg-floor}

A floor contact below refers to an interior local minimum of \(t\).
Thus \(\dd t=0\), \(\Hess t\ge0\), and
\(\Hess l/l=p\Hess t\ge0\).
The estimates hold at any such minimum, not only at the particular
height chosen in the continuation.

\begin{lemma}[Gauss Inequality at a floor contact]
\label{lem:fl-alg-floor-gauss}
For a symmetric tensor \(A\), define
\[
 \mathcal S(A)
 =\frac14(\tr_\perp A)^2-\frac12|A_\perp^{\mathrm{tf}}|^2
       +dA_{ee}\tr_\perp A-d|A_{e\perp}|^2.
\]
At an interior floor contact,
\begin{equation}
 \frac12e^{4t}R_{\hat g}
 \le \frac12R_g-v\Ric_g(e,e)+\mathcal S(H^f).
 \label{eq:fl-alg-floor-gauss}
\end{equation}
The expression \(v\Ric_g(e,e)\) is interpreted as \(\Ric_g(w,w)\);
all quantities have continuous, direction-independent values at
\(w=0\).
\end{lemma}

\begin{proof}
Consider the graph of \(f\) in the Riemannian warped product
\(g+l^2\dd s^2\). Its induced metric is \(\bar g\), its unit normal
has horizontal component \(-w\), and its vertical component has
length \(\sqrt d\).
Write \(T=l^{-1}\partial_s\) for the unit vertical vector.
The ambient curvature formulas are
\begin{align*}
 R_{\mathrm{amb}}&=R_g-2l^{-1}\Delta_gl,\\
 \Ric_{\mathrm{amb}}(X,Y)&=\Ric_g(X,Y)-l^{-1}\Hess l(X,Y),\\
 \Ric_{\mathrm{amb}}(T,T)&=-l^{-1}\Delta_gl,
\end{align*}
with zero mixed Ricci components.
Hence half the ambient term
\(R_{\mathrm{amb}}-2\Ric_{\mathrm{amb}}(n,n)\) in the graph Gauss
equation is
\[
 \frac12R_g-v\Ric_g(e,e)
                         -\frac vl\tr_\perp\Hess l.
\]
At the contact \(\dd l=0\); the covariant components of the graph
second form are therefore \(H^f\), up to the harmless common sign
coming from the choice of graph normal.
Tracing with the graph inverse metric \(A_d\) gives
\[
 \frac12\bigl((\tr_{A_d}A)^2-|A|_{A_d}^2\bigr)=\mathcal S(A).
\]
Indeed the axial square cancels, the mixed contribution is
\(-d|A_{e\perp}|^2\), and the transverse contribution is
\(\frac14(\tr_\perp A)^2-\frac12|A_\perp^{\mathrm{tf}}|^2\).
Finally, at \(\dd t=0\) the conformal scalar formula contributes
\(-4\tr(A_d\Hess t)\). We have thus proved the exact formula
\[
 \frac12e^{4t}R_{\hat g}
 =\frac12R_g-v\Ric_g(e,e)+\mathcal S(H^f)
       -\frac vl\tr_\perp\Hess l-4\tr(A_d\Hess t).
\]
Both final terms are nonpositive, proving
\eqref{eq:fl-alg-floor-gauss}.
\end{proof}

\begin{lemma}[Uniform floor coercivity]
\label{lem:fl-alg-floor-coercivity}
At a point with \(\dd t=0\), for an absolute constant \(C\),
\begin{equation}
 \mathcal T-\mathcal S(H^f)
 \ge \frac1{16}\mathcal T
        -C(1+N)^2\bigl(|K|^2+vF^2\bigr).
 \label{eq:fl-alg-floor-coercivity}
\end{equation}
In particular the coefficient \(1/16\) is independent of \(N\) and
of the graph gradient.
\end{lemma}

\begin{proof}
Subtracting \(\mathcal S(S)\) from
\eqref{eq:fl-alg-expanded-quadratic}, with \(\dd t=0\), gives
\begin{align}
 \mathcal T-\mathcal S(S)
  ={}&|P^{\mathrm{tf}}|^2+(1+d)|M|^2
       +\chi(1+d-\chi/2)q_1^2-dFq_1+\frac12F^2.
 \label{eq:fl-alg-floor-difference}
\end{align}
The identities
\[
 d\ge\chi,\qquad d-\chi=\frac{pv}{4}\chi,\qquad
 \frac{d^2}{\chi}=\frac{d(4+pv)}4\le\frac{N+4}{4}
\]
make its parameter dependence explicit.
For the part with cross coefficient \(-\chi\), Young's Inequality gives
\[
 \frac12F^2-\chi Fq_1+\chi(1+\chi/2)q_1^2
       \ge\frac18F^2+\frac56\chi q_1^2.
\]
The excess cross term satisfies
\[
 |(d-\chi)Fq_1|
 \le\frac13\chi q_1^2
           +\frac34\frac{(d-\chi)^2}{\chi}F^2
 \le\frac13\chi q_1^2+\frac3{64}p^2vF^2.
\]
It follows from \eqref{eq:fl-alg-floor-comparison} that
\begin{equation}
 \mathcal T-\mathcal S(S)
       \ge\frac18\mathcal T-\frac3{64}p^2vF^2.
 \label{eq:fl-alg-unshifted-floor}
\end{equation}

Because \(H^f=S-K\), expansion of the quadratic polynomial
\(\mathcal S\) gives
\[
 |\mathcal S(S-K)-\mathcal S(S)|
 \le C|K|\bigl(|P^{\mathrm{tf}}|+|M|+|F|+d|q_1|\bigr)+C|K|^2.
\]
Here we used \(\tr P=F-\chi q_1\) and \(\chi\le d\).
The estimate for \(d^2/\chi\) above and
\eqref{eq:fl-alg-floor-comparison} bound this by
\[
 \frac1{16}\mathcal T+C(1+N)|K|^2.
\]
Combining with \eqref{eq:fl-alg-unshifted-floor} proves the claim.
\end{proof}

\subsection{Differentiating the trace and the weighted tensor drift}
\label{sec:fl-alg-derivatives}

The next lemma records precisely the hypotheses on the continuation
right sides that enter the floor argument. They concern values and
first derivatives only. Let \(E\) be a fixed compact set containing
the support of \(\tr K\) and of all collar corrections. Fix
\(0<\delta<1\) and \(3/4<\gamma<1\), and positive smooth weights
\(\rho_0,\rho_1\) with
\[
 \rho_0\asymp r^{-3-\delta},\qquad
 \rho_1\asymp r^{-2\gamma}
\]
on the Euclidean end. Suppose
\[
 |K|\le Cr^{-2},\qquad |\nabla K|\le Cr^{-3}
\]
there. All constants defining \(E\), the weights, and these bounds
are fixed before \(N\).

For clarity denote the trace right side as a function of its inputs
by \(\mathfrak F(x,w,h,p)\), and its value on the current fields by
\(F\). The allowed forms are
\begin{align}
 \mathfrak F_1&=h+C(x), \notag\\
 \mathfrak F_{2,\theta}
    &=h+C(x)+\theta D_0(x,w,h), \qquad 0\le\theta\le1,\notag\\
 \mathfrak F_{3,\zeta}
    &=h+(1-\zeta)\{C(x)+D_0(x,w,h+\zeta b(x))\}
       +\zeta\{\tr_{A_\chi}K+D_1(x,w)\},
       \qquad 0\le\zeta\le1 .
 \label{eq:fl-alg-stage-right-sides}
\end{align}
Here \(C,b\) are fixed smooth functions, and
\(D_0,D_1\) are smooth bundle functions, supported over \(E\), that
extend smoothly to \(|w|\le1\).
Assume \(\partial_hD_0\ge0\).
On the bounded height interval under consideration, their values and
first derivatives in the base, fiber, and height variables are
bounded independently of \(N\).
Only \(A_\chi\), through \(p\), has \(N\)-dependent coefficients.
Base derivatives of bundle functions mean covariant derivatives with
their fiber arguments parallel transported.

\begin{lemma}[Floor derivative estimates]
\label{lem:fl-alg-floor-derivatives}
Assume the preceding hypotheses, \(h=\tau f\), and
\(|h|\le C_0\). At any point where \(\dd t=0\), for each fixed
\(\eta>0\) there are constants \(C_\eta\) and an integer \(m\),
independent of \(N\), the homotopy parameters, and the graph gradient,
such that
\begin{align}
 w(F)&\ge \tau a|\dd f|-\eta\mathcal T
              -C_\eta(1+N)^m(\rho_0+v\rho_1),
       \label{eq:fl-alg-trace-derivative}\\
 \left|\frac{\operatorname{div}(uA_\chi K(w,\cdot))}{u}\right|
      &\le\eta\mathcal T+C_\eta(\rho_0+v\rho_1).
       \label{eq:fl-alg-drift-derivative}
\end{align}
If in addition \(|h|\le Cr^{-\gamma}\) on the end, then
\begin{equation}
 |K|^2+vF^2\le C(\rho_0+v\rho_1).
 \label{eq:fl-alg-weighted-remainder}
\end{equation}
The constants may depend on the fixed data, weights, compact
corrections, height bound, and \(\eta\); they are uniform over the
three stages in \eqref{eq:fl-alg-stage-right-sides}.
\end{lemma}

\begin{proof}
\emph{The error weights.}
Since \(\delta+2\gamma<3\), the end decay implies
\[
 |\nabla K|\le C\sqrt{\rho_0\rho_1},\qquad
 |K|^2\le C\rho_0.
\]
Young's Inequality therefore gives
\begin{equation}
 a|\nabla K|\le C(\rho_0+v\rho_1).
 \label{eq:fl-alg-gradient-weight}
\end{equation}
Any bounded error supported in \(E\) is bounded by \(C\rho_0\),
because \(\rho_0\) has a positive minimum there.

\emph{The exact tensor-drift derivative.}
Put \(H=H^f\), \(k_0=K_{ee}\), \(m_0=K_{e\perp}\), and
\(P_H=H|_{e^\perp}\). At the point in question,
\[
 \nabla p=0,\qquad \nabla\log u=H(w,\cdot),
 \qquad \nabla w=H A_d.
\]
For \(c=(p+4)/(4+pv)\), Equation
\eqref{eq:fl-alg-smooth-matrix} reads
\[
 A_\chi K(w,\cdot)=K(w,\cdot)-cK(w,w)w.
\]
Its exact differentiated form is
\begin{align}
 \frac{\operatorname{div}(uA_\chi K(w,\cdot))}{u}
 ={}&a\{(\operatorname{div}K)(e)
                     -(1-\chi)(\nabla_eK)(e,e)\}\notag\\
 &+\chi(2\chi-1)H_{ee}k_0
      +2\chi\langle H_{e\perp},m_0\rangle
      +P_H:K_\perp-(1-\chi)(\tr P_H)k_0.
 \label{eq:fl-alg-exact-drift}
\end{align}
To see the cancellation explicitly, use a normal orthonormal frame
at the point. The entries of \(\nabla w=H A_d\) are
\[
 \nabla_e w_e=dH_{ee},\quad
 \nabla_e w_A=H_{eA},\quad
 \nabla_A w_e=dH_{Ae},\quad
 \nabla_A w_B=H_{AB}.
\]
The coefficient of \(H_{ee}k_0\), after adding
\(\langle\nabla\log u,A_\chi K(w,\cdot)\rangle\), is
\[
 d-3cvd-2(\partial_vc)v^2d+v\chi
 =\chi(2\chi-1).
\]
The mixed coefficient is \(1+d-2cv+v=2\chi\).
The remaining transverse terms are the last two terms of
\eqref{eq:fl-alg-exact-drift}; differentiating \(K\) produces its
first line. This proves the formula.

In particular every axial Hessian term in this divergence contains a
factor \(\chi\). Since \(H=S-K\) and
\(\tr P=F-\chi q_1\), the second line of
\eqref{eq:fl-alg-exact-drift} has absolute value at most
\[
 C|K|\bigl(|P^{\mathrm{tf}}|+|M|+|F|+\chi|q_1|\bigr)+C|K|^2
 \le\eta\mathcal T+C_\eta|K|^2.
\]
Here \(\chi|q_1|\le\sqrt{\chi}|q_1|\) and
\eqref{eq:fl-alg-floor-comparison} were used.
The first line is bounded by \(Ca|\nabla K|\).
Equation \eqref{eq:fl-alg-gradient-weight} proves
\eqref{eq:fl-alg-drift-derivative}, without an \(N\)-dependent constant.

\emph{The trace derivative in each stage.}
At fixed \((x,w,p)\), the height derivatives in
\eqref{eq:fl-alg-stage-right-sides} are respectively
\[
 1,\qquad 1+\theta\,\partial_hD_0,\qquad
 1+(1-\zeta)\partial_hD_0(x,w,h+\zeta b(x)),
\]
so each is at least one.
At \(\dd t=0\), the chain rule and \(h=\tau f\) give
\begin{equation}
 w(F)
 =\partial_h\mathfrak F\,\tau a|\dd f|
       +(\nabla_x\mathfrak F)(w)
       +D_w\mathfrak F(\nabla_w w),
 \qquad
 \nabla_w w=a(dH_{ee}e+H_{e\perp}).
 \label{eq:fl-alg-trace-chain}
\end{equation}
There is no \(\partial_p\mathfrak F\) term because \(\nabla p=0\).

The compact corrections have uniformly bounded first derivatives.
For the only noncompact correction, use
\[
 \tr_{A_\chi}K=\tr K-cK(w,w),\qquad
 D_wc[\xi]=2(\partial_vc)\langle w,\xi\rangle.
\]
The bounds
\[
 |c|\le C(1+N),\qquad
 |\partial_vc|\le C(1+N)^2
\]
give, on the stated height interval,
\[
 |\nabla_x\mathfrak F|
      \le C(1+N)^2\bigl(\mathbf 1_E+|\nabla K|\bigr),
 \qquad
 |D_w\mathfrak F|
      \le C(1+N)^2\bigl(\mathbf 1_E+|K|\bigr).
\]
The base derivative here also includes the bounded derivative of
\(b(x)\) inside \(D_0\).
The second term of \eqref{eq:fl-alg-trace-chain} is consequently
bounded below by
\(-C(1+N)^2(\rho_0+v\rho_1)\).
For its third term, the comparison
\eqref{eq:fl-alg-floor-comparison} and
\(d^2/\chi\le(N+4)/4\) yield
\[
 \begin{split}
 &C(1+N)^2(\mathbf 1_E+|K|)\,
       a\bigl(|H_{e\perp}|+d|H_{ee}|\bigr)\\
 &\hspace{12mm}\le
     \eta\mathcal T+
       C_\eta(1+N)^5(\mathbf 1_E+|K|^2).
 \end{split}
\]
For example, replace \(H\) by \(S-K\), use weighted Young Inequalities
against \(|M|^2+\chi q_1^2\), and absorb the remaining \(K\)-terms
into the same right side. Thus one may take \(m=5\), after enlarging
the constant. The weight observations at the start of the proof and
\(\partial_h\mathfrak F\ge1\) establish
\eqref{eq:fl-alg-trace-derivative}.

Finally, outside \(E\), all three right sides have the form
\[
 F=h+\zeta\tr_{A_\chi}K
\]
with \(0\le\zeta\le1\), where \(\zeta=0\) in the first two stages.
The eigenvalues of \(A_\chi\) lie in \((0,1]\), so
\(|F|\le |h|+C|K|\le Cr^{-\gamma}\) if the stated end height bound
holds. On \(E\), \(F\) is uniformly bounded by the compact hypotheses.
This proves \eqref{eq:fl-alg-weighted-remainder}.
\end{proof}

\begin{remark}
\label{rem:fl-alg-interfaces}
The preceding propositions and lemmas are identities and estimates
for smooth fields. Their application to the continuation uses the
right sides and compact corrections of
Proposition~\ref{prop:fl-homotopies}, the height bounds of
Lemma~\ref{lem:fl-height}, and the actual trace equation.
In particular no existence conclusion for the coupled system is
part of the scalar-curvature calculation.
\end{remark}


\begin{lemma}[Boundary flux and conformal mean curvature]
\label{lem:fl-boundary-identity}
On any face where \(f\) is constant, for the actual trace \(F\),
\begin{equation}
\label{eq:fl-boundary-identity}
 \frac{V_\nu}{u}
 =d(H+4\partial_\nu t)-H+w_\nu(F-P_B),\qquad
 H_{\hat g}=e^{-2t}\sqrt d\,(H+4\partial_\nu t).
\end{equation}
In particular the physical boundary condition gives
\(H_{\hat g}=-e^{-2t}\sqrt d\,N_B<0\).
\end{lemma}

\begin{proof}
The constant face value gives
\(\Hess f|_{TB}=(\partial_\nu f)\mathrm{II}_g\), and hence
\(\tr_{TB}H^f=w_\nu H\).
Using Equation~\eqref{eq:fl-alg-flux-identity}, or taking the normal
component of the original flux directly, gives
\[
 \frac{V_\nu}{u}
 =4d\partial_\nu t+\chi w_\nu q_1
 =4d\partial_\nu t+w_\nu(F-P_B)-vH.
\]
This is the first identity.  The normal in \(\bar g\) is
\(\sqrt d\,\nu\); its mean curvature is \(\sqrt d\,H\).
The latter follows also by subtracting the connection change from
the tangential Hessian above.  The conformal boundary formula for
\(\hat g=e^{4t}\bar g\) yields the second identity.
At \(\dd f=0\) it follows by smooth continuation, or by the same
direct computation with \(d=1\).
\end{proof}

\subsection{A specified four-stage homotopy}
\label{subsec:fl-homotopy}

\begin{proposition}[The continuation family]
\label{prop:fl-homotopies}
There is a continuous piecewise smooth family, with parameter
\(\xi\in[0,4]\), starting at Equations~\eqref{eq:fl-trace},
\eqref{eq:fl-scalar}, and \eqref{eq:fl-boundary}, and ending at the
zero trace solution and scalar equation with zero source and inner
flux.  In the first three stages the source is the \(\Xi\) just
defined, always evaluated using the current trace \(F\).
Every trace right side has derivative at least one in \(h\), at fixed
point, \(Z\), and \(\dd f\).  The outer values remain \(f=Z=0\).
\end{proposition}

\begin{proof}
Choose a smooth \(j:\mathbb R\to[0,1]\), zero for \(t\le0\) and one
for \(t\ge1\).  For \(0\le\lambda\le1\), define
\[
 V_\lambda=uA_\chi(4\nabla Z+\lambda K(w,\cdot)),\qquad
 B_K=A_\chi K(w,\cdot).
\]
Until the last scaling stage, the inner scalar condition is
\begin{equation}
\label{eq:fl-homotopy-one}
 \frac{(V_\lambda)_\nu}{u}
  =[1-(1-\lambda)j(t)]
       [\,\mathcal B_F-(1-\lambda)(B_K)_\nu\,].
\end{equation}

We specify the two fixed collar corrections.  Choose
\(\eta>0\) with \(6\eta<b_+-b_-\).  On black collars take a
directional cutoff \(\chi_-(w\cdot\nu_s)\), supported where
\(w\cdot\nu_s<-1/2\) and equal to one at \(-1\), and a nonincreasing
height cutoff \(\eta_-(h)\), equal to one for
\(h\le b_-+\eta\) and zero for \(h\ge b_-+2\eta\).
On white collars use \(\chi_+\), supported where
\(w\cdot\nu_s>1/2\) and equal to one at \(1\), and a nondecreasing
\(\eta_+\), zero for \(h\le b_+-2\eta\) and one for
\(h\ge b_+-\eta\).  Multiply by disjoint smooth collar cutoffs
\(\kappa_B\), equal to one on the faces, and set
\[
 D_0(x,w,h)=
 -A_0\sum_{\text{black collars}}\kappa_B\chi_-\eta_-
 +A_0\sum_{\text{white collars}}\kappa_B\chi_+\eta_+,
 \qquad A_0>2\sup_B\lvert H\rvert+1 .
\]
Thus \(\partial_hD_0\ge0\), \(D_0(x,0,h)=0\);
it is nonpositive at low heights \(h\le b_-+\eta\), nonnegative
at high heights \(h\ge b_+-\eta\), and zero in every compatible
black outward or white inward direction, for all lengths
\(0\le\lvert w\rvert\le1\).  All its derivatives on bounded height
ranges are independent of \(N\).

At the assigned face height, the fully corrected expression
\(F=h+C+D_0\), evaluated at the limiting normal directions
\(w=\pm\nu\), satisfies
\begin{equation}
\label{eq:fl-directional-inequalities}
 F(x,+\nu,b)\ge P_B+H,\qquad
 F(x,-\nu,b)\le P_B-H .
\end{equation}
These evaluations use \(\lvert w\rvert=1\) and \(\chi=0\) in the
smooth direction-free coefficients.  Equation~\eqref{eq:fl-compatible-values}
gives equality in the compatible direction.  In the opposite
direction the correction of magnitude \(A_0\) gives the inequality.

Take a smooth compactly supported extension \(b(x)\), constant at its
assigned face value on each smaller collar.  Put
\[
 D_1(x,w)=A_1\sum_B\kappa_B(x)\,w\cdot\nu_s,\qquad
 A_1>\sup_B\lvert H\rvert+1 .
\]
The collars are disjoint.  Thus \(D_1(x,0)=0\) and
\(D_1(x,\pm\nu)=\pm A_1\) on every face.
The family is as follows.
\begin{enumerate}[label=\textup{(\roman*)}]
 \item For \(0\le\xi\le1\), put \(\lambda=1-\xi\), retain
 \(F=h+C\) and \(h|_B=b\), and use
 \(\operatorname{div}V_\lambda=u\Xi\) with
 Equation~\eqref{eq:fl-homotopy-one}.
 \item For \(1\le\xi\le2\), put \(\lambda=0\), \(\alpha=\xi-1\),
 retain \(h|_B=b\), and set
 \(F=h+C+\alpha D_0(x,w,h)\).
 Keep the same scalar equation and boundary condition.
 \item For \(2\le\xi\le3\), put \(\lambda=0\), \(\zeta=\xi-2\),
 prescribe \(h|_B=(1-\zeta)b\), and set
 \begin{equation}
 \label{eq:fl-homotopy-three}
 F=h+(1-\zeta)\bigl[C+D_0(x,w,h+\zeta b(x))\bigr]
          +\zeta\bigl[\tr_{A_\chi}K+D_1(x,w)\bigr].
 \end{equation}
 Keep the same scalar equation and boundary condition.
 \item For \(3\le\xi\le4\), retain this trace equation at
 \(\zeta=1\), and multiply both \(\Xi\) and the right side of
 Equation~\eqref{eq:fl-homotopy-one} by \(s_\xi=4-\xi\).
\end{enumerate}
The prescriptions agree at their endpoints.  The height derivatives
are \(1\), \(1+\alpha\partial_hD_0\), and
\(1+(1-\zeta)\partial_hD_0\), respectively.  At a stage-(iii) face,
\(h+\zeta b=b\), so its limiting directional value is the convex
combination of the earlier corrected value and \(P_B+D_1(\pm\nu)\).
Equation~\eqref{eq:fl-directional-inequalities} therefore persists.

At \(\zeta=1\), the trace equation becomes
\(\tr_{A_\chi}H^f=h+D_1(x,w)\), with zero boundary values.
At a positive interior maximum of \(f\), \(w=0\), \(A_\chi=I\),
and \(D_1=0\), so \(l\Delta f=\tau f>0\), a contradiction.
A negative minimum is excluded similarly.  Hence \(f=0\) and \(t=Z\)
throughout the last stage.  At \(\xi=4\) the scalar equation is
\(\operatorname{div}(4L(Z)\nabla Z)=0\), with zero inner flux and
zero outer Dirichlet value.  Multiplication by \(Z\) gives \(Z=0\).
For a frozen-coefficient scalar solve at this endpoint, the output is
also identically zero for every input \(Z\).
\end{proof}

\subsection{Uniform height bounds and the tracefree tail}
\label{subsec:fl-height}

\begin{lemma}[Height and outer-boundary control]
\label{lem:fl-height}
Every smooth trace solution in Proposition~\ref{prop:fl-homotopies}
satisfies \(\lvert h\rvert\le C_0\), and its actual trace satisfies
\(\lvert F\rvert\le C_F\), independently of \(N,R,\xi\) and the input
\(Z\).  There are fixed \(r_0,C>0\) and a lower bound
\(R_{\mathrm{ht}}(N)\) such that, when
\(R\ge R_{\mathrm{ht}}(N)\),
\begin{equation}
\label{eq:fl-end-height}
 \lvert h\rvert\le C(r^{-\gamma}-R^{-\gamma})
       \quad(r_0\le r\le R),\qquad
 \lvert\dd f\rvert\le C\tau^{-1}R^{-1-\gamma}\quad\hbox{on }S_R .
\end{equation}
The constants in the bound are independent of \(N,R,\xi\).
For fixed \(N,\epsilon\), increasing the lower bound on \(R\)
ensures \(t>-\epsilon\) at the outer boundary of every such solution.
\end{lemma}

\begin{proof}
The corrections \(D_0,D_1\) vanish at \(w=0\).  At an interior
extremum of \(f\), the first two stages give
\(\tr K+l\Delta f=h+C\), and the third gives
\(\tr K+l\Delta f=h+(1-\zeta)C+\zeta\tr K\).
The maximum and minimum inequalities bound \(h\) by
\(\lVert\tr K\rVert_\infty+\lVert C\rVert_\infty\) and the
boundary heights.  Boundedness of the corrections and
\(\lvert\tr_{A_\chi}K\rvert\le3\lvert K\rvert\) then bounds \(F\).
No floor has been used.

Choose \(r_0\) outside the supports of \(C,D_0,D_1,b\) and
\(\tr K\).  The last stage has \(h=0\).  In the other stages the
tail equation takes the form
\begin{equation}
\label{eq:fl-tail-trace}
 \frac{l}{\tau\sqrt D}\tr_{A_\chi}\Hess h-h
       +\theta\,\tr_{A_\chi}K=0,\qquad 0\le\theta\le1,
\end{equation}
where \(\theta=1\) in the first two stages and \(1-\zeta\) in the
third.  Since \(\tr K=0\),
\[
 \tr_{A_\chi}K=(\chi-1)K(e,e),\qquad
 1-\chi=\frac{(p+4)v}{4+pv}.
\]
For the radial test \(\psi=r^{-\gamma}\), with its radial gradient
axis, Equation~\eqref{eq:fl-data-decay} gives uniformly in
\(0<\chi\le1\)
\[
 \tr_{A_\chi}\Hess\psi
 =\gamma\bigl((\gamma+1)\chi-2\bigr)r^{-2-\gamma}
                 +O(r^{-3-\gamma})
 \le-c_\gamma r^{-2-\gamma}
\]
after increasing \(r_0\).

Consider \(\phi=C(r^{-\gamma}-R^{-\gamma})\), choosing \(C\)
large enough to dominate \(C_0\) at \(r=r_0\) for
\(R\ge2r_0\).  In the part \(r\le R/2\), its height is at least
\(C(1-2^{-\gamma})r^{-\gamma}\); since \(\gamma<1<2\),
this dominates the bounded adverse tensor term \(O(r^{-2})\)
after fixing \(r_0\) sufficiently large.

In the part \(r\ge R/2\), evaluate coefficients at a possible
comparison contact, so that the given input \(Z\) and the test
gradient agree with the solution's.  Put
\(z=l\lvert\dd\phi\rvert/\tau\), noting
\(\lvert\dd\phi\rvert\asymp C r^{-1-\gamma}\).
If \(z\ge1\), the magnitude of the negative Hessian term is at
least \(c r^{-1}\), whereas the adverse tensor term is at most
\(C_Kr^{-2}\).  If \(z\le1\), the former is at least
\(c(l/\tau)C r^{-2-\gamma}\), while the latter is at most
\[
 C_K(N+4)(l/\tau)^2 C^2 r^{-4-2\gamma}.
\]
Their ratio is at most \(C(N+4)z/r\le C(N+4)/r\).
Thus a lower bound \(R_{\mathrm{ht}}(N)\), enlarged after
\(r_0,C\) have been fixed, makes \(\phi\) a strict upper barrier.
Its negative is a lower barrier by the same absolute estimates.
Height monotonicity allows comparison with the actual input \(Z\).
This proves the first bound in Equation~\eqref{eq:fl-end-height}.
Tangential derivatives vanish on \(S_R\); comparison at their common
outer value gives the stated normal derivative bound.

Finally, \(Z=0\) on \(S_R\), and the map \(t\mapsto Z(t,\dd f)\)
is increasing.  Its value at \(t=-\epsilon\) is negative whenever
\[
 \ell^2\lvert\dd f\rvert^2<e^{8\epsilon}-1.
\]
The boundary slope bound makes this true once
\[
 C^2\ell^2\tau^{-2}R^{-2-2\gamma}<e^{8\epsilon}-1 .
\]
This is a permissible further lower bound on \(R\) at fixed
\(N,\epsilon\).
\end{proof}

\subsection{Exclusion of every floor contact}
\label{subsec:fl-floor}

We now fix the constants in the source.  Only estimates at a proposed
floor contact are needed; no upper bound on \(Z\) or on \(\lvert\dd f\rvert\)
is assumed.

\begin{lemma}[Error ledger at an interior floor contact]
\label{lem:fl-floor-ledger}
There exist \(c_0>0\) and polynomial constants \(C_N\) such that
every smooth first-three-stage solution with
\(t\ge-\epsilon\), at an interior point with \(t=-\epsilon\),
satisfies
\begin{equation}
\label{eq:fl-floor-divergence}
 u^{-1}\operatorname{div}V_\lambda
 \ge c_0\mathcal T+\tau a\sigma-C_N(\rho_0+v\rho_1).
\end{equation}
Here \(c_0\) can be chosen independently of \(N,R,\xi\).
\end{lemma}

\begin{proof}
At the contact \(\dd t=0\), \(\Hess t\ge0\), \(p=N\), and
\(\dd p=0\).
Equation~\eqref{eq:fl-alg-identity} and
Lemma~\ref{lem:fl-alg-floor-gauss} give
\[
 \frac{\operatorname{div}V}{u}
 \ge 8\pi\mu-\tfrac12R_g+8\pi J(w)+v\Ric(e,e)
        +\mathcal T-\mathcal S(H^f)+w(F)-F\tr K .
\]
The constraint gives exactly
\[
 8\pi\mu-\tfrac12R_g
       =\tfrac12\bigl((\tr K)^2-\lvert K\rvert^2\bigr).
\]
Lemma~\ref{lem:fl-alg-floor-coercivity} retains
\(\tfrac1{16}\mathcal T\), with an error
\(C_N(\lvert K\rvert^2+vF^2)\).
The first-derivative estimates for the current trace and the omitted
flux, from Lemma~\ref{lem:fl-alg-floor-derivatives}, retain \(\tau a\sigma\)
and each cost an arbitrarily small fixed fraction of \(\mathcal T\).
They have polynomial remainders.

For clarity, every coefficient remainder has the following global
control, with constants independent of \(N,R,\xi\):
\begin{align*}
 \lvert K\rvert^2&\le C\rho_0,&
 vF^2&\le Cv\rho_1,&
 \lvert F\tr K\rvert&\le C\rho_0,\\
 v\lvert\Ric\rvert&\le Cv\rho_1,&
 a(\lvert\nabla K\rvert+\lvert J\rvert)
      &\le C(\rho_0+v\rho_1).
\end{align*}
The first bound uses \(\delta<1\); the second follows from
Lemma~\ref{lem:fl-height} and the tail expression
\(F=h+\zeta\tr_{A_\chi}K\).
The trace term has compact support.  The curvature bound follows
from \(\Ric=O(r^{-3})\).
Finally, on the end
\(\lvert\nabla K\rvert+\lvert J\rvert=O(r^{-3})\) and
\[
 \lvert\nabla K\rvert+\lvert J\rvert
       \le C\sqrt{\rho_0\rho_1},
 \qquad \delta+2\gamma<3 .
\]
Weighted Young's Inequality gives the last displayed estimate.
All compactly supported coefficient derivatives are bounded by a
constant times \(\rho_0\), which has positive minimum on their
fixed compact support.

Choose the two small derivative fractions so that their sum is less
than \(1/32\), and subtract
\((1-\lambda)\operatorname{div}(uA_\chi K(w,\cdot))/u\).
This proves Equation~\eqref{eq:fl-floor-divergence}, after renaming
the positive universal coefficient.  No fixed-\(N\) regularity
constant enters this ledger.
\end{proof}

\begin{proposition}[Floor separation along the full homotopy]
\label{prop:fl-floor}
Fix \(0<\delta_0<\min\{c_0/2,1/4\}\), where \(c_0\) is from
Lemma~\ref{lem:fl-floor-ledger}.  There is a polynomial choice of
\(\Pi_N\) such that, for all sufficiently large \(N\), followed by
all \(R\ge R_{\mathrm{floor}}(N,\epsilon)\), no smooth solution
of Proposition~\ref{prop:fl-homotopies} satisfies \(t\ge-\epsilon\)
and attains \(t=-\epsilon\).  Every such solution has \(t>-\epsilon\)
on the closed truncation.
\end{proposition}

\begin{proof}
Enlarge \(R_{\mathrm{floor}}\) to include the outer-boundary
requirement of Lemma~\ref{lem:fl-height}.  Thus an outer contact is
impossible.

At an inner contact in one of the first three stages, \(j(t)=0\).
Equation~\eqref{eq:fl-homotopy-one} cancels the omitted drift on
both sides and gives \(V_\nu/u=\mathcal B_F\).
Lemma~\ref{lem:fl-boundary-identity} implies
\[
 \partial_\nu t=-\tfrac14(N_B+H)<0 .
\]
This contradicts the nonnegative inward derivative at a boundary
minimum.

At an interior contact, Equation~\eqref{eq:fl-floor-divergence}
has right side at least
\(2\delta_0\mathcal T+\tau a\sigma-C_NW_N\).
Choose the polynomial \(\Pi_N\) so large that
\[
 (\Pi_N-C_N)W_N\ge 2\rho+\rho_0 .
\]
This is possible since \(W_N\ge\rho_0\) and
\(\rho/\rho_0\) is bounded.  At the floor \(m_0=1\);
comparison with the prescribed source would then give
\[
 0\ge
 \delta_0\mathcal T+(1-\delta_0)\tau a\sigma
           +(\Pi_N-C_N)W_N-\rho
 \ge\rho+\rho_0>0,
\]
a contradiction.

In the last stage \(f=0\), \(t=Z\), \(u=L\), and \(v=0\).
At an interior floor minimum,
\(\mathcal T=\frac12(\lvert K\rvert^2+(\tr K)^2)\).
Enlarge the same polynomial so that
\(\delta_0\mathcal T+\rho-\Pi_N\rho_0<0\) at every such point.
For \(s_\xi>0\), the scalar equation then has
left side \(4L\Delta t\ge0\) and strictly negative right side.
At an inner floor minimum its boundary condition gives
\(4\partial_\nu t=s_\xi(-H-N_B)<0\).
Both are impossible.  If \(s_\xi=0\), multiplication of
\(\operatorname{div}(4L(t)\nabla t)=0\) by \(t\), with zero outer
value and zero inner flux, gives \(t=0\).
\end{proof}

The estimates in this section concern every smooth candidate solution
and its homotopy, and thus supply the boundary exclusion needed for
continuation.  Existence is established in
Proposition~\ref{prop:fc-existence}, using the global estimates and
regularity arguments of the following sections.
